% SPDX-License-Identifier: GPL-3.0-or-later
% Copyright (C) 2026 Alessandro Greco
\chapter{流量关联、指纹识别与 Sybil 攻击}\label{cap:26-attacchi}
\aggiornatoal{2026 年 9 月 23 日}{本章引用的研究论文；tor 0.4.9.13；实验网络}

\begin{obiettivi}
学完本章后，你将能够：
\begin{itemize}[nosep]
  \item 根据攻击者的观察位置对针对 Tor 匿名性的攻击分类；
  \item 解释端到端流量关联，并在实验网络中测量自己产生的流量（实验~16）；
  \item 描述网站指纹识别以及相关评估的局限性；
  \item 识别 Sybil 攻击及 Tor 的防御机制，从入口守卫到带宽测量。
\end{itemize}
\end{obiettivi}

\Cref{cap:04-modello-minaccia} 的威胁模型已经表明：Tor 无法抵御同时观察通信两端的攻击者。二十多年的研究量化了这一警告，并发现实际可用的观察点比最初想象的更多，也更容易接近。本章从防御者的视角分析这些攻击：它们需要什么条件、能够取得什么效果，以及我们已经研究过的哪些协议机制可以提高攻击成本。

% ---------------------------------------------------------------------------
\section{攻击者的位置}\label{sec:att-posizioni}

按照攻击者能够观察到的内容，可以将攻击分为：
\begin{description}
  \item[一个端点] 互联网接入服务商、本地网络或审查者可以观察客户端到入口守卫的连接；攻击者也可能观察出口到目的地的连接。在客户端与入口守卫之间，它能观察加密流量的时间、数据量和方向；出口与目的地之间是否有应用层加密则取决于协议。网站指纹识别可以利用客户端侧的元数据，而无需读取内容。
  \item[两个端点] 攻击者通过在入口和出口同时运营中继，或者控制二者路径上的网络，观察通信两端，从而进行流量关联。
  \item[网络内部] 恶意中继，包括中间中继，可以测量、延迟、标记或丢弃 cell。
  \item[底层互联网] 中继连接经过的自治系统（AS）和互联网交换节点（IXP）等底层网络基础设施。
\end{description}

% ---------------------------------------------------------------------------
\section{端到端流量关联}\label{sec:att-correlazione}

为了保持可用性，Tor 有意避免对 cell 施加混合网络式的长时间延迟或重新排序（\cref{sec:radici-latenza}）。因此，流量的\emph{形状}——每个时间区间内通过多少字节——通常仍可在电路两端被识别，但会受到延迟、缓冲、多路复用和其他流量的影响（\cref{fig:correlazione}）。同时观察两端的人，不需要破解任何加密，就可以比较这些形状并判断两条流是否对应。

\begin{figure}[htbp]
\centering
% SPDX-License-Identifier: GPL-3.0-or-later
% Copyright (C) 2026 Alessandro Greco
%% EPUB 无障碍替代文本。
\altfig{端到端流量关联的单向合成示例，并非实测数据。上方图表示在客户端到守卫链路上每 250 毫秒观察到的加密流量相对大小；下方是出口到服务器链路的相对流量。两条独立设定的示例序列有相似但不完全相同的突发模式。图中下方序列延后一段 250 毫秒，峰值高度和噪声也发生变化。加密不一定隐藏时间和流量规模特征；缓冲、协议开销、多路复用和并发流量会改变模式。这种固定延迟并非 Tor 的普遍特性，统计关联也不能保证判断正确。}
% 合成曲线，非真实观测
\begin{tikzpicture}[font=\sffamily\scriptsize,x=0.32cm,y=0.09cm]
  \def\serieIngresso{0/6,1/0,2/0,3/9,4/12,5/2,6/0,7/0,8/0,9/5,10/14,11/14,12/3,13/0,14/0,15/7,16/0,17/0,18/11,19/4,20/0,21/0,22/8,23/2}
  \def\serieUscita{0/0,1/5,2/0,3/0,4/8,5/11,6/3,7/0,8/0,9/1,10/4,11/13,12/12,13/4,14/0,15/0,16/6,17/0,18/1,19/10,20/5,21/0,22/1,23/7,24/3}
  \foreach \x/\h in \serieIngresso {
    \fill[torviola!70] (\x,20) rectangle ++(0.8,\h);
  }
  \foreach \x/\h in \serieUscita {
    \fill[torverde!70!black] (\x,0) rectangle ++(0.8,\h);
  }
  \draw[->] (-0.5,20) -- (27,20) node[right] {$t$};
  \draw[->] (-0.5,0) -- (27,0) node[right] {$t$};
  \node[left,align=right] at (-0.6,26) {入口\\观察者\\（客户端至守卫）};
  \node[left,align=right] at (-0.6,6) {出口\\观察者\\（出口至服务器）};
  \draw[<->,torgrigio] (10,17) -- node[above,font=\sffamily\tiny] {示例：偏移一个区间} (11,17);
  \node[text=torgrigio] at (13,40) {每 250\,ms 区间相对流量};
  \node[text=torgrigio,font=\sffamily\tiny] at (13,36.5) {合成示例，非真实测量};
\end{tikzpicture}

\caption{单向流量关联的两条\emph{合成}示例曲线，并非真实测量：分别表示客户端到入口守卫、出口到服务器链路的相对流量。突发模式相似但并不完全相同，图中的一个 250\,ms 区间偏移仅为示例，而非 Tor 的固定延迟。加密可以保护内容，但未必隐藏流量时间和规模特征。缓冲、开销、多路复用和并发流量可能改变两端模式；统计相似不能单独证明流量关联成立。}
\label{fig:correlazione}
\end{figure}

\paragraph{匹配有多容易？}2018 年，DeepCorr 训练神经网络来识别 Tor 流量对。在每条流约有 900 个数据包的条件下，它达到了 96\,\% 的准确率；相同条件下，之前的统计方法只有 4\,\%~\cite{nasr2018deepcorr}。对攻击者而言，最困难的部分通常不是比较流量，而是占据适当的观察位置。

\paragraph{攻击者同时观察到两端的概率有多高？}Johnson 等人使用 Tor 实际的路径选择算法，对典型用户进行了持续数月的模拟~\cite{johnson2013users}。假设攻击者运行具有中等带宽的恶意中继，在研究所建模的不同用户类型中，有 80\,\% 的用户在六个月内至少发生过一次去匿名化。面对单个自治系统，部分常见位置的用户在三个月内暴露概率接近 100\,\%。这些历史模拟结果并不意味着每条连接都遭到监视；真正的含义是攻击机会随时间积累。正是这一推理推动了更长寿命的入口守卫和每位客户端使用更少守卫的设计（\cref{sec:guard}）。

\paragraph{以互联网基础设施为攻击者。}RAPTOR 识别出自治系统增加观察机会的三种途径~\cite{sun2015raptor}：利用互联网路径的不对称性，因为在通信两端各观察\emph{一个方向}的流量（包括 TCP 确认报文）就可能足够；利用随时间自然发生的路由变化；以及通过 BGP 劫持或拦截地址前缀。这些技术都无须修改 Tor。防御方向包括监控互联网路由，以及可能采用考虑网络拓扑的路径选择。

\paragraph{主动攻击。}控制通信一端的攻击者，可以故意通过延迟或集中发送数据包，为流量\emph{刻入}可识别的节奏，并在另一端寻找这一模式。\cref{sec:tor1-debolezze} 中讨论的加密层标记攻击，是这种思想的极端例子。早在 2005 年，Murdoch 和 Danezis 就表明，攻击者并不总需要直接观察真实链路。恶意服务器可以为发往受害者的流量施加特定节奏；攻击者再测量自己通过候选中继所建立电路的延迟，如果某些中继的延迟也随该节奏变化，就说明它们可能位于受害者电路上~\cite{murdoch2005lowcost}。

% ---------------------------------------------------------------------------
\section{实验：关联自己产生的流量}\label{sec:lab-attacchi}

\begin{laboratorio}[16]{端到端流量关联}
\textbf{目标。}在实验网络中证明，流量模式能够穿过电路而保留下来，并用于识别某条出口侧流量属于哪个客户端。\\
\textbf{前置条件。}实验~11；实验网络正在运行。所有流量都由实验脚本自行生成。\\
\textbf{预计时间。}1 小时。
\end{laboratorio}

\file{labs/zh-Hans/实验室/correlation.py} 脚本在容器内的端口 8095 启动服务器，并启动六个客户端。各客户端使用不同的 SOCKS 凭据，在 \opt{IsolateSOCKSAuth} 启用时，不同客户端的流不能共用同一电路（\cref{sec:lab-stem}）。在 20 秒内，客户端随机交替发送 2\,KB 数据块和保持静默。两个\enquote{观察者}分别记录数据：
\begin{itemize}
  \item 入口观察者记录每个客户端何时写入数据。这里简化了实际入口守卫链路观察者所能看到的内容：真实数据还会被封装进 cell 和 TLS，但关联所需的时间模式会穿过这两层。
  \item 出口观察者就是我们运行的服务器。它记录各条连接上的字节到达时间，同时收到一个用于核对结果的四字节标识符，但匹配算法不使用该标识符。
\end{itemize}
两个观察者都按 250\,ms 的时间区间汇总字节。对每个客户端，我们计算其时间序列与全部六条出口侧流量之间的 Pearson 相关系数。相关性最高的流被选为可能的对应项。脚本实际上知道正确答案，因为客户端在最初四个字节中提供了身份信息，但这一信息只用于评估匹配结果，不用于选择最相关的流。

\begin{lstlisting}[language=bash,numbers=none]
$ docker exec -w /tmp torbook-network python3 /lab/correlation.py 1
客户端 0： +0.79 -0.11 +0.12 +0.17 +0.11 -0.01  -> 出口流 0 匹配
客户端 1： -0.11 +0.80 -0.02 +0.07 +0.18 +0.08  -> 出口流 1 匹配
客户端 4： +0.13 +0.02 +1.00 +0.21 +0.10 +0.21  -> 出口流 2 匹配
客户端 3： +0.21 +0.09 +0.33 +0.81 -0.07 +0.03  -> 出口流 3 匹配
客户端 2： +0.17 +0.05 +0.10 -0.09 +0.85 -0.03  -> 出口流 4 匹配
客户端 5： -0.00 +0.01 +0.23 +0.06 -0.13 +0.92  -> 出口流 5 匹配
正确匹配数：6/6
\end{lstlisting}
列对应按建立连接的到达顺序排列的出口侧流量，因此其顺序与客户端编号并不一致。每行的正确匹配都很明显，相关系数在 0.79 到 1.00 之间，其余通常接近零。这是一次参考运行的节选，而不是精确预测：实际运行时流的顺序与系数可能不同，也可能只匹配成功 5/6。这里既不需要解密密钥，也不需要查看解密后的 cell，只需要数据量和时间信息。实验网络规模很小，负载也较轻，因此延迟极小。DeepCorr 的历史研究表明，在所研究的条件下，更复杂的分类器甚至可以利用真实网络中几百个数据包进行匹配。此实验包含在本书的实验验证中。

% ---------------------------------------------------------------------------
\section{网站指纹识别}\label{sec:att-wf}

如果攻击者只能观察客户端到入口守卫的链路，就不能执行端到端关联，但仍可能\emph{识别}正在访问的网站。每个网页都会在双向连接上产生具有一定特征的 cell 序列，包括数量、顺序和集中发送的模式。攻击者先通过 Tor 访问目标网站，记录这些序列并训练分类器，再将受害者的流量与已学习的模式比较。

相关研究一直是分类器与防御机制之间的持续较量：
\begin{itemize}
  \item 2011 年，Panchenko 等人利用流量大小、时间和方向特征，在包含 775 个网站的封闭集合中，将识别率从 3\,\% 提升到 55\,\%。他们还展示了有效的\emph{开放世界}攻击，即受害者可能访问训练集之外的网站~\cite{panchenko2011website}。
  \item 2018 年，Deep Fingerprinting 使用卷积神经网络，在 95 个未受保护的网站上取得超过 98\,\% 的准确率；面对曾被考虑用于 Tor 的自适应填充防御 WTF-PAD，准确率仍达到 90\,\%~\cite{sirinam2018deep}。
\end{itemize}

\begin{approfondimento}[忽视基础发生率的谬误]
极高的准确率数字需要谨慎解释。Juarez 等人指出，许多研究采用了有利于攻击者的假设：用户一次只访问一个网页、受害者与攻击者使用相同版本的 Tor Browser，而且训练数据在攻击发生前不久收集~\cite{juarez2014critical}。更根本的问题是，在开放世界中，受监视的网站只占所有访问的一小部分。即使分类器的误报率很低，在数百万次无关访问中产生的误报也可能多于正确识别。这与针对罕见疾病的筛查测试面临的统计陷阱相同；真实攻击者同样无法回避。
\end{approfondimento}

\paragraph{防御机制。}通道填充（\cref{cap:07-canali}）使空闲时段更难识别，却并非专门为抵御网页指纹攻击而设计。tor 还实现了更一般的\emph{电路填充}框架：状态机可以按计划发送伪装用 cell。在本书参考日期，已经部署的状态机可以帮助洋葱服务引介与汇合电路看起来更像普通电路（\cref{cap:14-onion-services}，\proposta{302}）。能在可接受带宽开销下广泛有效防御网站指纹识别的方法，仍是尚待解决的研究问题。

% ---------------------------------------------------------------------------
\section{Sybil：大量伪造中继身份}\label{sec:att-sybil}

\textbf{Sybil 攻击}是指向网络引入大量身份。在 Tor 中，攻击者添加大量中继，希望它们频繁被选为入口守卫和出口中继。三类防御会提高攻击成本：
\begin{description}
  \item[测量带宽] 中继如果虚报可用带宽，就可能吸引过多客户端。Bauer 等人在 2007 年展示了这种基于描述符的欺骗手段~\cite{bauer2007lowresource}。如今，共识带宽权重以目录权威节点的测量为基础，而不是仅信任中继的声明（\cref{sec:dir-voto}）；因此，吸引大量流量通常需要真正提供带宽。
  \item[路径约束] 属于同一家族或同一个 IPv4 /16 网络的两台中继，不能同时出现在一条电路上（\cref{sec:percorso-vincoli}）。诚实运营者会声明家族；不诚实的运营者至少仍受到 IP 地址分布的约束。
  \item[网络监控] 权威节点可以排除中继或为其设置 \texttt{BadExit} 标志。志愿者与研究人员调查具有共同特征的中继群组，并监控出口行为。例如，2014 年一项出口扫描发现了 25 台操纵流量或配置不当的中继，其中包括 HTTPS 拦截、DNS 篡改与 \emph{sslstrip}~\cite{winter2014spoiled}。
\end{description}
\Cref{cap:27-casi-rete} 将研究此后发现的恶意中继群组，并说明为什么监控必须持续进行。

\paragraph{针对洋葱服务的 Sybil 攻击。}在 v2 中，负责某个服务的 HSDir 位置可以预测。攻击者因此能够生成特定中继身份，让它们占据这些位置。Biryukov 等人在 2013 年利用这种手法，测量服务热度、使服务不可访问，或协助定位攻击~\cite{biryukov2013trawling}。在 v3 中，位置取决于盲化密钥和每天更新的共享随机值（\cref{cap:14-onion-services}），因此显著增加了提前预计算和定向布置恶意 HSDir 的难度；但当期随机值公布后仍可算出负责的 HSDir，并不能完全消除恶意目录节点。

% ---------------------------------------------------------------------------
\section{发现入口守卫}\label{sec:att-guard}

入口守卫保护客户端，却也成为一个有吸引力的攻击目标：一旦知道它是谁，攻击者就知道应该在哪里观察。面对洋葱服务，攻击者可以通过不断请求连接等方式触发大量电路建立，直到某个恶意中继被选为紧邻守卫的中间中继。2006 年发表的早期同类攻击，也推动了入口守卫机制的发展~\cite{overlier2006locating}。Vanguards（\cref{sec:vanguards}）则额外引入一层在一段时间内保持固定的中间中继。这会拖慢攻击者逐层追踪至真正入口守卫的过程，但并不保证守卫无法被发现。

% ---------------------------------------------------------------------------
\section*{本章小结}
\begin{itemize}
  \item 攻击的区别在于观察位置：一个端点、两个端点、网络内部，或底层互联网路由。
  \item 流量关联无需破解加密，只需比较两端的流量形状。攻击难点在于占据适当位置，而观察机会会随时间积累。
  \item 网站指纹识别仅根据一端的加密流量识别网页。评估时必须考虑基础发生率和现实条件。
  \item Sybil 防御包括带宽测量、家族声明、/16 约束，以及对出口和中继群组的监控。在 v3 中，提前操纵 HSDir 位置更难，但当期负责的 HSDir 仍可计算。
  \item 入口守卫和 Vanguards 使发现第一跳的攻击更加缓慢且昂贵。
\end{itemize}

% ---------------------------------------------------------------------------
\section*{练习}

\begin{esercizio}[1]\label{es:26-posizione}
针对本章每种攻击，确定攻击者至少需要哪种观察位置，并找出能够提高该攻击成本的 Tor 防御机制。
\end{esercizio}

\begin{esercizio}[1]\label{es:26-base}
某分类器以 0.95 的概率正确识别受监视的网站，并以 0.001 的概率把无关访问误报为该网站。如果每 100\,000 次访问中只有一次真正访问了被监视的网站，那么所有警报中有多少比例是正确的？
\end{esercizio}

\begin{esercizio}[2]\label{es:26-bin}
重复实验~16，先把时间区间改为两秒，再把实验总时长缩短为五秒，最后让所有客户端采用相同的流量节奏。哪种修改对匹配结果的破坏最大？为什么？
\end{esercizio}

\begin{esercizio}[2]\label{es:26-difesa}
修改实验~16 的客户端，让它们以固定速率发送数据块，即使没有实际数据也发送填充数据。这会消耗多少额外带宽？流量关联结果会怎样变化？
\end{esercizio}

\begin{esercizio}[3]\label{es:26-probabilita}
使用 \file{path\_selection.py} 和实验~15 的真实共识，计算如下概率：攻击者控制 5\,\% 的守卫带宽和 5\,\% 的出口带宽，在六个月内至少一次同时占据某客户端电路的两端。比较固定入口守卫与每条电路都重新选择守卫的情形。请明确假设六个月内的电路数量，并仅在教学模型中按照带宽权重进行独立抽样；解释该模型为何不能预测真实网络。
\end{esercizio}

% ---------------------------------------------------------------------------
\section*{延伸阅读}
\begin{itemize}
  \item A.~Johnson 等，\titolo{Users Get Routed}~\cite{johnson2013users}。
  \item Y.~Sun 等，\titolo{RAPTOR: Routing Attacks on Privacy in Tor}~\cite{sun2015raptor}。
  \item M.~Juarez 等，\titolo{A Critical Evaluation of Website Fingerprinting Attacks}~\cite{juarez2014critical}。
  \item \emph{Anonymity Bibliography}：\url{https://www.freehaven.net/anonbib/}，可用于检索后续研究。
\end{itemize}
